% !TeX root = ../../Book2.tex
\section{边界正则性}
\label{b2:sec:1603}

内部差商可以向任意坐标方向平移，靠近边界时却可能越界。对于零 Dirichlet 条件，平坦边界允许沿切向作差商，而法向二阶导数则由方程补回。曲边界先化成平坦边界；这个坐标变换消耗的正则性，决定了边界应满足的假设。

\subsection{边界平坦化}
\label{b2:subsec:160301}

称有界区域为 \(C^{1,1}\) 区域，是指边界在有限个局部正交坐标中可写成 \(x_d=\psi(x')\)，其中 \(\psi\) 一阶可微且梯度 Lipschitz 连续，区域位于图像的一侧。局部使用
\[x=\Phi(y)=(y',y_d+\psi(y')),\qquad y_d>0\]
把半球或半盒映到区域内部。映射及其逆一阶导数有界，二阶弱导数也有界；因此它们保留局部 \(H^1\)、\(H^2\) 正则性，并给出相应范数的双向控制。

\ProseParagraphBreak

若 \(v(y)=u(\Phi(y))\)、\(\zeta(y)=\varphi(\Phi(y))\)，链式法则和换元把型写成
\[
 \int A(x)\nabla_xu\cdot\nabla_x\varphi\,\mathrm dx
 =\int\widetilde A(y)\nabla_yv\cdot\nabla_y\zeta\,\mathrm dy,
\]
其中
\[\widetilde A=|\det D\Phi|\,D\Phi^{-1}(A\circ\Phi)D\Phi^{-\mathsf T},
 \qquad\widetilde f=(f\circ\Phi)|\det D\Phi|.\]
图像平移的行列式为 \(1\)，保留一般写法便于检查其他坐标。由于 \(D\Phi^{-\mathsf T}\xi\) 的长度与 \(|\xi|\) 双向可比，\(\widetilde A\) 仍一致椭圆；若 \(A\in W^{1,\infty}\)、\(\Phi\in C^{1,1}\)，则 \(\widetilde A\in W^{1,\infty}\)。右端 \(\widetilde f\) 仍属于 \(L^2\)。这正是下一步差商证明需要的三个条件。

\ProseParagraphBreak

只有 Lipschitz 边界时，\(D\Phi\) 一般只有本性有界性，变换后的系数未必具有有界一阶导数。此时切向差商中的系数误差便不再受控。边界光滑性不是为了画图方便，而是通过坐标系数直接进入估计。

\subsection{Dirichlet 问题}
\label{b2:subsec:160302}

先考虑半球 \(B_2^+=B(0,2)\cap\{x_d>0\}\)。设 \(u\in H^1(B_2^+)\) 在平坦部分的迹为零，且弱解 \(-\operatorname{div}(A\nabla u)=f\)。测试函数可以支撑到平坦部分，但须在那里取零迹，并在球面部分附近消失。这样的测试等式由原 Dirichlet 弱形式经局部化得到。

\begin{lemma}[边界切向估计]
\label{b2:lem:16-tangential-estimate}
设 \(A\in W^{1,\infty}(B_2^+)\) 有界且一致椭圆，\(f\in L^2(B_2^+)\)，\(u\) 如上具有零平坦边界迹。则对 \(k<d\) 及 \(j\le d\)，\(\partial_k\partial_j u\in L^2(B_1^+)\)，且
\[\sum_{k<d,j\le d}\lVert\partial_k\partial_j u\rVert_{L^2(B_1^+)}
 \le C\bigl(\lVert f\rVert_{L^2(B_2^+)}+\lVert u\rVert_{H^1(B_2^+)}\bigr).\]
\end{lemma}
\begin{proof}
取在 \(B_1\) 上为 \(1\)、支集在 \(B_{3/2}\) 内的光滑截断 \(\eta\)。对切向方向 \(e_k\)，平移不改变 \(x_d\)，所以 \(D_h^ku\) 的平坦边界迹仍为零。于是 \(-D_{-h}^k(\eta^2D_h^ku)\) 是合法零迹测试函数；它在曲面边界附近为零，没有额外边界项。

在半球上重复内部估计的离散分部积分。所有平移都沿切向，故积分换元不越过平坦边界；截断又避免了球面边界。主项仍有椭圆下界，系数差商仍由 \(\lVert\nabla A\rVert_\infty\) 控制，右端仍由差商的 \(H^1\) 界控制。因此 \(\eta\nabla D_h^ku\) 的 \(L^2\) 范数独立于 \(h\) 有界。差商引理在半球内部测试并取极限，得到所列切向及混合导数及范数界。
\end{proof}

现在只缺 \(\partial_{dd}u\)。在分布意义展开方程，得
\[
 a_{dd}\partial_{dd}u
 =-f-\sum_{(i,j)\ne(d,d)}a_{ij}\partial_{ij}u
       -\sum_{i,j}(\partial_i a_{ij})\partial_j u.
\]
右端各项均属于 \(L^2(B_1^+)\)：所有非纯法向二阶项已由引理控制，系数导数有界，解的一阶导数平方可积。又 \(a_{dd}\ge\lambda>0\)，故除以它得到 \(\partial_{dd}u\in L^2\)。展开与除法均可在分布中合法进行，因为 \(a_{dd}\) 及其倒数是局部 Lipschitz 乘子。至此得到局部边界 \(H^2\) 估计。

\ProseParagraphBreak

这一步没有沿法向把 \(u\) 随意平移，也没有先把 \(u\) 作二阶可微延拓。零迹保证切向测试合法，椭圆性保证纯法向最高阶项的系数可逆，两者承担不同职责。一维时没有切向方向，最后的方程直接恢复二阶导数，恰与一维计算一致。

\subsection{边界 \texorpdfstring{\(H^2\)}{H2} 估计}
\label{b2:subsec:160303}

\begin{theorem}[全局 Dirichlet 二阶正则性]
\label{b2:thm:16-boundary-h2}
设 \(\Omega\subset\mathbb R^d\) 是有界 \(C^{1,1}\) 区域，\(A\in W^{1,\infty}(\Omega)\) 有界且一致椭圆，\(f\in L^2(\Omega)\)。若 \(u\in H_0^1(\Omega)\) 弱解 \(-\operatorname{div}(A\nabla u)=f\)，则 \(u\in H^2(\Omega)\)，且
\[\lVert u\rVert_{H^2(\Omega)}\le C\bigl(\lVert f\rVert_2+\lVert u\rVert_{H^1(\Omega)}\bigr)\le C'\lVert f\rVert_2.\]
常数依赖区域的有限边界坐标、系数上界与 Lipschitz 范数及椭圆下界。
\end{theorem}
\begin{proof}
覆盖 \(\partial\Omega\) 的有限个 \(C^{1,1}\) 坐标把方程变成上一小节的半球模型。坐标变换保留零迹、\(L^2\) 右端及 Lipschitz 椭圆系数；切向估计与恢复纯法向导数给出每个较小边界坐标内的 \(H^2\) 界。余下的紧内部区域由有限个球覆盖，应用内部二阶估计。将有限个局部积分相加，并使用坐标变换的双向范数控制，得到第一项不等式。上一章的能量估计与 Poincaré 不等式控制 \(\lVert u\rVert_{H^1}\le C\lVert f\rVert_2\)，得到第二项。
\end{proof}

对非齐次边界，若给定延拓 \(g\in H^2(\Omega)\)，令 \(w=u-g\)。右端变为 \(f+\operatorname{div}(A\nabla g)\in L^2\)，故
\[\lVert u\rVert_{H^2}\le C\bigl(\lVert f\rVert_2+\lVert g\rVert_{H^2}\bigr).\]
这里比弱存在性所需的 \(g\in H^1\) 多一阶正则性。若只有粗糙边界数据，内部解仍可能光滑，但上述全局二阶结论不再由本定理保证。

\ProseParagraphBreak

对足够光滑的区域和数据，可以在边界坐标中继续切向微分，再用方程恢复高阶法向导数。每次操作都要控制新的系数和坐标导数，因此需要同步提高边界光滑性。全局高阶定理不能仅由内部自举一句带过；本章只完整证明后续最常用的 \(H^2\) 层次。

\subsection{边界光滑性的作用}
\label{b2:subsec:160304}

取平面扇形 \(0<r<1\)、\(0<\theta<\omega\)，其中 \(\pi<\omega<2\pi\)，令 \(\alpha=\pi/\omega\in(1/2,1)\)。函数
\[v(r,\theta)=r^\alpha\sin(\alpha\theta)\]
在扇形内部调和，并在两条直边上为零。它的梯度量级为 \(r^{\alpha-1}\)，故近角点的梯度平方积分由 \(\int_0^\varepsilon r^{2\alpha-1}\,\mathrm dr\) 控制而有限；二阶导数量级为 \(r^{\alpha-2}\)，相应积分 \(\int_0^\varepsilon r^{2\alpha-3}\,\mathrm dr\) 发散。

\ProseParagraphBreak

选径向光滑截断 \(\chi\)，使它在角点附近为 \(1\)，在外圆边界附近为零，令 \(u=\chi v\)。则 \(u\in H_0^1\)，但 \(u\notin H^2\)；源项 \(f=-\Delta u\) 只支撑在离开角点的环带内，因而光滑且平方可积。局部化并在内圆 \(r=\varepsilon\) 上积分，边界余项按 \(\varepsilon^\alpha\) 消失，验证 \(-\Delta u=f\) 确实弱成立。于是即便主部是常系数、源项很好，凹角仍可破坏全局 \(H^2\) 正则性。

\ProseParagraphBreak

这个区域具有 Lipschitz 边界，却在凹角处不属于 \(C^{1,1}\)。例子不意味着所有非 \(C^{1,1}\) 区域都失败，例如某些凸多边形仍有更好的结论；它说明仅凭任意 Lipschitz 边界不能保证本节的全局二阶估计。内部估计在离角点有正距离的每个子域上仍完全成立，两个结论并不冲突。
